\section{Introduction}
\label{sec:intro}

%% Web relevance must be explicit on the first page (CFP requirement); target track: Web Infrastructure and Agentic Systems.
LLM agents now reach the Web through tool protocols. Wealth-management assistants, travel agents and CRM copilots all expose REST services as tools, increasingly through the Model Context Protocol (MCP)~\cite{anthropic2024mcp}, and let a language model decide which tools to call and with which arguments~\cite{schick2023toolformer,yao2023react}. How these services are \emph{described} to the agent and how the resulting calls are \emph{executed} are therefore questions of Web infrastructure. Together they decide the latency, cost and reliability of every agent built on these services.

Tool calls are rarely independent. Answering ``how much does the steady portfolio I hold charge?'' requires listing the customer's portfolios, deciding which one the customer means, reading its fee-schedule identifier, and only then querying the fee schedule. In a ReAct loop~\cite{yao2023react}, each dependent step costs one LLM round, and every round re-reads all previous tool outputs. On the multi-tool split of FinMCP-Bench~\cite{zhu2026finmcp}, reference trajectories take \todo{5.7} LLM rounds, and the median task returns about \todo{38k} characters of tool output.

The strongest current answer is to let the agent \emph{write code}. CodeAct~\cite{wang2024codeact} and programmatic tool calling (PTC)~\cite{anthropic2025advancedtooluse} have the model write one program that chains the calls in a sandbox, keeping intermediate results out of the context. Plan-then-execute compilers~\cite{xu2023rewoo,kim2024llmcompiler} do the same with a plan language. These approaches remove an LLM round only when the step is \emph{mechanical}: a field lookup, a filter on a numeric field, a computation. Two obstacles remain.

\paragraph{Semantic decisions return control to the model.} At some steps the next argument or branch depends on a judgment over data that only exists at run time: ``the steady one'' among twelve returned portfolios, ``the announcements that could explain the drop'', whether a product's risk text matches the customer's concern. A program cannot make this judgment. It can hard-code a brittle string match, or it can return the data to the model and pay for another round. Our screening of FinMCP-Bench (\S\ref{sec:screening}) finds such decisions in most multi-tool tasks.

\paragraph{Code over unfamiliar tools is written blind.} MCP tools usually declare inputs but not outputs, cross-tool dependencies or business rules. A model writing code against them must guess field names, or first print a result and look at it, which again adds rounds. Nothing checks the program before it runs, so invalid calls and wrong arguments surface only as run-time failures.

We present \jif, a \emph{graph-guided tool compiler} that addresses both obstacles. The agent writes the tool chain once, in a restricted Python subset (\pyif), against typed stubs that the compiler derives from a domain graph. Every semantic decision is written as a call to a System-1 primitive, such as \texttt{if jev.noul(\ldots):} or \texttt{p = jev.choice(\ldots)}. This is where the name comes from: a System-1 model is the evaluator of the program's \texttt{if}s. The compiler type-checks the program against the graph, turns each primitive into a well-posed question using the graph's entity semantics, batches and parallelizes, and lowers the program to an executable flow. At run time, tools are executed by the runtime and decisions are evaluated by Jev~\cite{typesafe2026jev}, a calibrated non-generative model. A confidence gate returns a decision to the LLM only when it is uncertain.

Two designs make this work.
\begin{itemize}
  \item \textbf{A split object--operation graph (\sog) as the compiler's input.} The object graph models entities, properties, relations and business rules. The operation graph models tool-backed operations, and its data-flow edges are derived from property production and consumption, so nobody annotates them by hand. Existing graph-guided agents use a graph to \emph{select or plan} tools~\cite{lee2026innout,navi2025,gap2025}. \jif uses the \sog to \emph{compile}. The operation graph supplies types and data-flow checks; the object graph supplies the labels, meanings and rules that make each run-time decision well posed.
  \item \textbf{\jif: System-1 models as the evaluators of program conditionals.} Choosing among a few typed candidates, checking a yes/no guard, or keeping items in a filter needs neither reasoning traces nor generated text. It needs a fast closed-set judgment with an honest probability. Placing these judgments \emph{inside} the program as typed primitives is what lets a chain with semantic steps run end to end. Conformal calibration bounds the probability that the program acts on a wrong decision.
\end{itemize}

The result is fast and cheap by construction: the LLM is called to plan and to answer, and otherwise only when a decision is uncertain. It does not trade away accuracy. Static checking removes invalid calls before execution, graph semantics reduce wrong-entity decisions, and provenance checks reject fabricated arguments.

\paragraph{Contributions.}
\begin{itemize}
  \item \textbf{\jif, a graph-guided tool compiler} for agent tool chains: the \pyif language with System-1 primitives, graph-based static checking, elaboration of decisions with entity semantics, batching and parallelization, and lowering to an executable flow with confidence-gated fallback (\S\ref{sec:compiler}).
  \item \textbf{Split object--operation graphs} as the compiler's input, with derived data flow, object-scoped disclosure, rendering of typed stubs, and a construction procedure from tool schemas and call logs (\S\ref{sec:sog}).
  \item \textbf{An evaluation} on FinMCP-Bench, screened into tasks with and without semantic decisions, together with an anonymized advisory deployment and synthetic scaling suites. A factorial study on top of programmatic tool calling (graph guidance $\times$ \jif) and targeted ablations isolate each design (\S\ref{sec:experiments}). \todo{release code, graphs and screened splits}
\end{itemize}
